\chapter{Discussion}
\label{ch:discussion}

This chapter interprets the results of Chapter~\ref{ch:results} against the research questions of
Section~\ref{sec:contributions}, relates them to prior work, and consolidates the threats to validity.

\section{What the Internal Readout Adds}
\label{sec:d-internal}

The results answer RQ1 and RQ2 sharply. At the output (Table~\ref{tab:infill},
Figure~\ref{fig:infill}), the four unlearning methods are indistinguishable. Each drives forget accuracy to
the never-knew floor, and no output number ranks them. The internal readout
(Table~\ref{tab:ranks}, \ref{tab:midstack}; Figures~\ref{fig:logit8b}, \ref{fig:lenstable}) resolves the
single output verdict into four different mechanisms and shows that none of them removes the fact's
internal trace under the readouts applied.
The clearest single demonstration is the confidence contradiction of Figure~\ref{fig:confidence}. Three
unlearned models emit a wrong token at essentially the same near-certain confidence as the model that
knows the right one. Any metric built on output probability is therefore blind by construction to the
distinction the thesis cares about, because the quantity it reads, output confidence, is genuinely
uninformative about internal knowledge here.

The deeper point is that suppression and erasure are \emph{different states of the model} that happen to
produce the same behaviour on the queried distribution. The relearning and quantization attacks of
Section~\ref{sec:r-attacks} make this concrete. The same output-forgotten fact is one fine-tuning step, or
one change of numerical precision, away from returning, and how far away depends on the internal state that
the output cannot see. An evaluation that stops at the output is thus not merely incomplete but
systematically optimistic, since it counts every suppression as an erasure.

\section{The Four Mechanisms}
\label{sec:d-mechanisms}

Reading the mid-stack against the output (Table~\ref{tab:midstack}) and confirming it with the calibrated
tuned lens (Table~\ref{tab:tuned}) yields a coherent mechanistic account of the four methods, which the
robustness experiments then corroborate.

\emph{IDK} is the purest surface case. It changes the output policy (``I don't know'') without disturbing
the representation, so the birthdate token remains recoverable through to the output (rank $9$), and only
free-generation infill hides it. Consistent with an intact representation, IDK resists relearning yet leaks the most under
quantization. \emph{NPO} and \emph{GradDiff} share a late-suppression signature. They keep the token
accessible mid-stack and push it down only in the final layers, a late operation on an otherwise intact
computation; GradDiff pushes hardest at the output. \emph{RMU} is the one method that acts early and deep,
scrambling the representation at its intervention layers so that the token is buried mid-stack as well as at
the output; but precisely because this is a large-magnitude perturbation that a few relearning steps
reverse rather than a removal, RMU snaps back fastest under relearning. The unifying observation is that every method leaves the token recoverable
somewhere, at rank $147$ even for RMU under the calibrated lens, so on this benchmark unlearning masks
knowledge rather than deleting it.

\section{Why \texorpdfstring{$D_\text{mech}$}{D\_mech} is Faithful}
\label{sec:d-faithful}

RQ3 is answered by the meta-evaluation (Table~\ref{tab:faithful}). On Hubble-8B, $D_\text{mech}$ is the only
metric, mechanistic or lexical, that is both strongly discriminative and substantially paraphrase-stable,
with a composite faithfulness score far above any lexical metric and a ranking that holds
under every aggregation we tried (Section~\ref{sec:r-crossbench}). The mechanism of this result deserves
emphasis because it is
not simply ``read internals''. The na\"ive internal metric, the minimum rank over the band, \emph{fails}
invariance as badly as the lexical metrics (shift $1.000$). A meaning-preserving rewording shifts the
absolute internal rank just as it shifts the output probability. What makes $D_\text{mech}$ paraphrase-stable
is its \emph{relative} construction. Because it subtracts the trajectory's distance to the never-knew anchor
from its distance to the knows anchor, a rewording that moves the whole trajectory by a roughly constant
amount moves both distances together, and the subtraction cancels most of the common-mode shift while
preserving the sign of the knowledge signal. This stability is thus a designed property of the anchor
difference, not a lucky feature of hidden states, which is why the distance construction, and not the mere
use of internals,
is the contribution. This is also why $D_\text{mech}$ is built on the parameter-free logit lens. The
cancellation requires a single fixed transport shared by the fact and both anchors, and only the logit lens
provides one without introducing trained or estimated components into the metric itself.

\section{The Three Lenses as Metric Bases}
\label{sec:d-lenses}

RQ4 concerns the choice of lens. The three lenses form a ladder of increasing fidelity to the model's own
computation, from the correlational logit lens through the calibrated tuned lens to the interventionally-
grounded Jacobian lens, but higher fidelity does not make a better \emph{metric} basis. The tuned lens, trained to reproduce
the output, ``skips ahead'': Figure~\ref{fig:lenscompare} shows it reporting the answer at rank ${\approx}10$
from the first layer for \emph{every} model, including the suppressed ones, so it collapses the very
trajectory differences $D_\text{mech}$ reads. It is an excellent output predictor and a poor knowledge
probe, which is exactly why it serves here as a cross-check (confirming the logit-lens ordering,
Table~\ref{tab:tuned}) rather than as the metric base. The Jacobian lens is a first-order,
interventionally-grounded readout and, once corrected, reproduces the taxonomy with a smooth, well-ordered
readout (Figure~\ref{fig:jacobian}); it even distinguishes a lingering correlational trace from a
sensitivity-based signal, keeping NPO's fact buried at every layer where the logit lens still detects a
trace. But its dependence on a reference corpus, its first-order approximation, and its weaker behaviour at
1B make it a confirmatory readout rather than a metric base, and we do not read it as proof of causal
necessity. The
logit lens wins for the metric not because it is the most faithful readout but because it is the only
\emph{reproducible, parameter-free} one, and $D_\text{mech}$ restricts itself to the band where it is
reliable.

\section{Scale Dependence}
\label{sec:d-scale}

The cross-scale results (Table~\ref{tab:ranks}, Figure~\ref{fig:logit1b}) carry a methodological warning.
The same method can internalise unlearning at one scale and not another: GradDiff and RMU bury the fact at
8B but leave it accessible at 1B, and only NPO is robust across both. This interacts with the metric. At 1B,
where the base model memorises the facts only weakly, the knows and never-knew anchors are closer together,
so $D_\text{mech}$'s separation is correspondingly weaker. The practical implication is that unlearning
conclusions drawn at a single scale, and especially from a single output metric, can invert across scale, as
they do for IDK, which the output labels ``forgetting'' at 1B while the logit lens shows the fact merely
output-suppressed. Faithful evaluation therefore benefits from strong memorisation (a well-separated knows
anchor), which the 8B setting provides and the 1B setting only partially does.

\section{Relation to Prior Work}
\label{sec:d-related}

These findings extend three lines of prior work. The robustness literature has repeatedly shown that
unlearning can be undone by quantization \citep{zhang2024quantization}, relearning
\citep{lynch2024eight,fan2025relearn}, and black-box or weight probing \citep{thaker2024blackbox,
deeb2024weights}; our contribution is to show, on a controlled benchmark with ground-truth references,
\emph{where inside the model} the recoverable knowledge sits and to summarise it in a single faithful
metric. Closest to our internal readout, \citet{patil2023candelete} use logit-lens attacks to recover
edited facts from hidden states; we differ in targeting the modern unlearning methods rather than knowledge
editing, in using the lens as a systematic readout rather than an attack, and, crucially, in distilling it
into a metric validated by the field's own faithfulness meta-evaluation. Finally, where mechanistic
interpretability has met unlearning it has been used to \emph{perform} unlearning more robustly, by
localising and editing the responsible components \citep{guo2024mechanistic,sae2025clamping}; we use it
instead to \emph{measure} unlearning, filling the evaluation-side gap that OpenUnlearning
\citep{dorna2025openunlearning} identifies but addresses only for output metrics, and taking up the
interpretability-and-metrics agenda that Hubble \citep{allegro2025hubble} explicitly invites.

\section{Threats to Validity}
\label{sec:d-threats}

We restate the limitations consolidated. The faithfulness meta-evaluation is run at 8B with $n=44$
birthdate facts and a \emph{single} fixed paraphrase and wrong-target realisation, so its variance is across
facts, not across paraphrase styles; the multi-paraphrase robustness check
(Section~\ref{sec:r-invariance}) supports the ordering but a design that samples many rewordings would
harden it. The faithfulness advantage is scale-dependent and weaker at 1B (Section~\ref{sec:d-scale}). The
facts are treated as the unit of analysis, but they share model checkpoints, prompt templates, and
overlapping leave-one-out anchors, so the inferential statistics are benchmark-conditional rather than
estimates over an unrestricted population of facts. The primary Hubble analysis uses a single attribute
(birthdates); Section~\ref{sec:r-crossbench} replicates the meta-evaluation on OpenUnlearning TOFU
(author facts with multi-token answers) at three scales to test how far the result travels. The composite
faithfulness score is a convenience for ranking; the individual invariance and discrimination effect sizes
are the primary evidence, and both the within-pool Wilcoxon $p$-values and the TOST equivalence test are
uninformative at this sample size (they saturate), which is why we rely on effect sizes rather than on the
equivalence bound. The per-fact relearning route to
validating faithfulness returned a null (Section~\ref{sec:r-attacks}), so no claim rests on per-fact
prediction. Finally, one negative pool (the Perturbed model scored against a wrong birthdate) leaks slightly
on $D_\text{mech}$, because a wrong date still traces a date-like path. Reporting the two negative pools
separately makes this explicit: on Hubble-8B the discrimination effect size is a perfect $|\delta|=1.000$
against the clean never-knew pool but $0.653$ against the wrong-target pool (combined $0.827$), so the
positive-negative gap is carried mainly by the never-knew pool. This is noted but does not affect the
paraphrase-stability conclusion, which is where the metric is decided.
